% Standalone circuit figure for the IEEE Access manuscript.
% Compile once with pdflatex to produce fig_ecm_circuit.pdf, then include it in the paper.
\documentclass[border=2pt]{standalone}
\usepackage{amsmath}
\usepackage{times}
\usepackage[american]{circuitikz}
\begin{document}
\begin{circuitikz}[american voltages, font=\footnotesize, line width=0.7pt]
\coordinate (srcB) at (0,0);
\coordinate (srcT) at (0,1.6);
\coordinate (r0L)  at (0.45,1.6);
\coordinate (r0R)  at (1.75,1.6);
\coordinate (rcL)  at (2.45,1.6);
\coordinate (rcR)  at (4.45,1.6);
\coordinate (pos)  at (5.15,1.6);
\coordinate (neg)  at (5.15,0);
\draw (srcB) to[V,invert,l_={$\mathrm{OCV}(z)$}] (srcT);
\draw (srcT) -- (r0L) to[R,l={$R_0$}] (r0R) -- (rcL);
\draw (rcL) -- (2.45,2.20) to[R,l={$R_1$}] (4.45,2.20) -- (rcR);
\draw (rcL) -- (2.45,1.00) to[C] (4.45,1.00) -- (rcR);
\node[font=\footnotesize] at (3.85,1.18) {$C_1$};
\draw (rcR) -- (pos) (pos) to[open,v^={$\hat V(t)$}] (neg) (neg) -- (srcB);
\draw[-latex,thick] (4.60,1.85) -- node[above=1mm] {$I(t)$} (5.05,1.85);
\fill (pos) circle (1.3pt);
\fill (neg) circle (1.3pt);
\end{circuitikz}
\end{document}
